
\textbf{Proposition 6.2.2.} \uline{Pour toute section positive} \(\underline{n}\) \uline{de} \(\mathbb Z\)\uline{, le faisceau étale en groupes} \(\underline{\Der}(\Gl(\underline{n}))\) \uline{des commutateurs de} \(\Gl(\underline{n})\) \uline{est représentable par}
\(\operatorname{Sl}(\underline{n})=\Ker(\det:\Gl(\underline{n})\to{G}_m)\).

\textbf{Corollaire 6.2.3.} \uline{Toute représentation}
\begin{equation}\tag{6.2.4}
 \alpha:G=\Gl(\underline{n}_1)\times\cdots\times\Gl(\underline{n}_r)\longrightarrow{G}_m
\end{equation}
\uline{est de la forme}
\begin{equation}\tag{6.2.5}
 \alpha(Y)(s)=(\det(s_1))^{\underline{m}_1}\cdots(\det(s_r))^{\underline{m}_r},
\end{equation}
\uline{où} \(Y\in\operatorname{Ob}(\Sch/S)\), \(s=(s_1,\ldots,s_r)\in G(Y)\) \uline{et} \(\underline{m}_1,\ldots,\underline{m}_r\) \uline{sont des sections de} \(\mathbb Z\) \uline{uniquement déterminés par} \(\alpha\).

\textbf{Corollaire 6.2.6.} \uline{Avec les notations de} 6.1. \uline{on a}
\[
 \det\bigl(T_\alpha(\underline{E}_1,\ldots,\underline{E}_r)\bigr)
 =
 \det(\underline{E}_1)^{\otimes \underline{m}_1}\otimes\cdots\otimes
 \det(\underline{E}_r)^{\otimes \underline{m}_r},
\]
\uline{où} \(\underline{m}_1,\ldots,\underline{m}_r\) \uline{sont des sections de} \(\mathbb Z\) \uline{uniquement déterminées par} \(\alpha\).

\textbf{Remarque 6.2.7.} On peut déterminer les exposants \(\underline{m}_1,\ldots,\underline{m}_r\) dans (6.2.5) et (6.2.6) en substituant successivement dans (6.2.5)
\(s_i=\lambda.\ \operatorname{id}.(\lambda\in{G}_m(Y))\), \(s_j=\operatorname{id}_{\Gl(\underline{n}_j)(Y)}\), pour \(j\ne i\) \((i,j=1,\ldots,r)\), d'où
\[
 \det(s_i)=\lambda^{\underline{n}_i}\quad\text{et}\quad
 \alpha(Y)(s)=\lambda^{\underline{n}_im_i}.
\]
Par exemple on déduit ainsi de (6.2.6)
\begin{align}
 \det(\underline{E}_1\otimes \underline{E}_2)&=(\det(\underline{E}_1))^{\otimes \underline{n}_2}\otimes(\det(\underline{E}_2))^{\otimes \underline{n}_1}, \tag{6.2.8}\\
 \det(\underline{E}_1\oplus \underline{E}_2)&=(\det(\underline{E}_1))\otimes(\det(\underline{E}_2)) \tag{6.2.9}
\end{align}
et
\begin{equation}\tag{6.2.10}
 \det\bigl(\bigwedge\limits^{\underline{m}} \underline{E}_1\bigr)=(\det(\underline{E}_1))^{\binom{\underline{n}_1-1}{\underline{m}-1}}
\end{equation}
etc.

